% =========================================================
\section{Prueba integrada del sistema}
\label{sec:prueba_integrada}
% =========================================================

Las pruebas de la \cref{sec:validacion} ejercitan cada subsistema por
separado y sin instrumentos. Esta sección examina el comportamiento del
sistema completo en sesiones reales de laboratorio, con el láser
emitiendo, el osciloscopio adquiriendo y el módulo de temperatura
operando de forma simultánea.

Conviene precisar de entrada qué se pretende demostrar. Las sesiones no
caracterizan el agua: el arreglo no normaliza por la energía del pulso,
según se estableció en la \cref{sec:equipos}, y la señal registrada
lleva impresa la respuesta del transductor. Lo que se evalúa es la
herramienta. Si capturas tomadas bajo los mismos parámetros producen
resultados equivalentes dentro de una dispersión acotada, si un cambio
deliberado en un parámetro del láser se distingue de esa dispersión, y
si el registro de sesión conserva la información necesaria para afirmar
ambas cosas después sin recurrir a la memoria del operador, entonces el
sistema entrega lo que declara entregar.

La evidencia proviene de dos grupos de sesiones. El principal son las
sesiones del 3 y del 5 de noviembre de 2026, tomadas con la versión
congelada del código y con el formato de registro de 34 columnas
(\cref{sec:formato_datos}); con ellas se evalúan la repetibilidad, la
estabilidad del eje de tiempos y el comportamiento del ruido en ambos
modos automáticos. El segundo es la sesión del 25 de agosto de 2026,
anterior a la congelación, que se conserva porque es la única que
contiene un cambio deliberado de un parámetro del láser.

% ---------------------------------------------------------
\subsection{Sesiones con la versión congelada}
\label{sec:sesiones_congeladas}
% ---------------------------------------------------------

Las sesiones se realizaron sobre agua en enfriamiento pasivo, con el
osciloscopio a \SI{2}{\milli\volt} por división y
\SI{200}{\micro\second} por división, registros de \num{25000}
muestras y el láser en disparo continuo. En todas ellas el
M.~en~C. Arturo Ronquillo Arvizu operó el sistema en el laboratorio y
el autor lo supervisó a distancia. La \cref{tab:sesiones_congeladas}
resume las capturas automáticas consideradas.

\begin{table}[htbp]
  \centering
  \caption{Capturas automáticas de las sesiones con la versión
    congelada. $N$: adquisiciones promediadas por captura. El CV es el
    coeficiente de variación del valor pico a pico; el CV residual es el
    de los residuos tras ajustar una recta del valor pico a pico contra
    la temperatura registrada.}
  \label{tab:sesiones_congeladas}
  \small
  \setlength{\tabcolsep}{4pt}
  \begin{tabular}{l l c c c c c}
    \toprule
    Sesión & Modo & Capturas & $N$ & CV & CV residual & $R^2$ \\
    \midrule
    \texttt{20261103\_194433} & tiempo & 8 & 100 &
      \SI{2.1}{\percent} & \SI{1.4}{\percent} & 0.62 \\
    \texttt{20261103\_200010} & temperatura & 6 & 416--452 &
      \SI{2.3}{\percent} & \SI{0.7}{\percent} & 0.92 \\
    \texttt{20261105\_191232} & temperatura, $\pm\SI{0.1}{\celsius}$ & 5 &
      728--924 & \SI{2.1}{\percent} & \SI{1.8}{\percent} & 0.41 \\
    \texttt{20261105\_200231} & temperatura, $\pm\SI{0.02}{\celsius}$ & 6 &
      14--140 & \SI{6.2}{\percent} & \SI{3.8}{\percent} & 0.70 \\
    \texttt{20261105\_201925} & tiempo & 10 & 100 &
      \SI{4.4}{\percent} & \SI{3.8}{\percent} & 0.33 \\
    \bottomrule
  \end{tabular}
\end{table}

No todas las capturas de esas fechas entran en la tabla, y los
criterios de exclusión forman parte del resultado. La sesión
\texttt{20261103\_200010} se limita a sus capturas \texttt{m0002} a
\texttt{m0007}: de la \texttt{m0008} en adelante el contador del
instrumento registró más adquisiciones de las que el láser pudo
producir, es decir, el promedio arrastraba las capturas anteriores. Son
las seis capturas con exceso que el testigo de entonces aceptó
(falla~15, \cref{tab:fallas_software}). En sentido opuesto, las
capturas \texttt{m0006} y \texttt{m0007} de esa sesión y siete de la
\texttt{20261103\_194433} llevan la marca de error del testigo anterior
y se incluyen, porque su conteo de adquisiciones es coherente con el
ritmo del láser: son falsos positivos de la misma falla. La captura
\texttt{m0010} de la \texttt{20261103\_194433} se excluye porque
promedió 84 de las 100 adquisiciones solicitadas.

La sesión \texttt{20261105\_184959} se descartó completa. El contador
de adquisiciones no se reinició entre ventanas (de 966 a 3304) y los
ritmos aparentes, de 14 a 59 adquisiciones por segundo, son imposibles
con el láser a unos \SI{10}{\hertz}. La causa probable es que se
presionó el botón de marcha y paro del panel frontal, que deja la
condición de paro del instrumento en un estado distinto al que el
sistema programa. A diferencia del 3 de noviembre, el testigo corregido
marcó todas las capturas de esa sesión en el momento de tomarlas. A
partir de la siguiente sesión, el sistema reafirma el modo de
adquisición y la condición de paro antes de cada orden de marcha.

\subsubsection{Repetibilidad}
\label{sec:repetibilidad}

La métrica empleada para la amplitud es el valor pico a pico calculado
sobre la traza completa, una vez restada la línea base. No se utiliza
una ventana temporal fija, porque el origen del eje de tiempos lo fija
el operador con el control de posición horizontal del osciloscopio y no
es reproducible entre sesiones, según se discute en la
\cref{sec:origen_tiempos}.

El coeficiente de variación dentro de cada sesión va de
\SI{2.1}{\percent} a \SI{6.2}{\percent}. Esa cifra mezcla dos cosas que
conviene separar. En el modo por temperatura la amplitud desciende con
la temperatura de la muestra; en la sesión \texttt{20261103\_200010},
la única con ventanas de duración regular y sin capturas anómalas, la
recta contra la temperatura explica el \SI{92}{\percent} de la varianza
y la dispersión residual es de \SI{0.7}{\percent}. En las demás, la
temperatura explica una fracción menor y la dispersión residual se
mantiene entre \SI{1.4}{\percent} y \SI{3.8}{\percent}.

Las dos sesiones en modo por tiempo muestran el mismo patrón. La
amplitud cae durante las primeras cuatro o cinco capturas
---\SI{5.7}{\percent} el 3 de noviembre y \SI{12.8}{\percent} el 5 de
noviembre--- y después se estabiliza: en el tramo estable el
coeficiente de variación es de \SI{0.25}{\percent} y
\SI{0.34}{\percent}, respectivamente. En régimen estable, por lo tanto,
la herramienta repite al nivel de unas décimas de por ciento; la mayor
parte de la dispersión de la \cref{tab:sesiones_congeladas} es una
deriva al inicio de la secuencia. Su causa no puede atribuirse con los
datos disponibles: el arreglo no mide la energía del pulso, de modo que
una variación de la emisión del láser y un cambio de la muestra
producen en el registro el mismo efecto.

La deriva entre sesiones es mayor que la deriva dentro de una sesión. A
temperaturas cercanas a \SI{34}{\celsius}, el valor pico a pico fue de
\SI{9.85}{\milli\volt}, \SI{8.39}{\milli\volt} y
\SI{6.92}{\milli\volt} en tres sesiones del mismo día, una diferencia de
hasta \SI{42}{\percent}. Las amplitudes registradas por el sistema solo
son comparables dentro de una misma sesión.

\subsubsection{Estabilidad del eje de tiempos}

Con el eje reconstruido mediante la \cref{eq:conv_tiempo}, el instante
del máximo de la amplitud absoluta se mantuvo entre
\SI{66.64}{\micro\second} y \SI{66.72}{\micro\second} en las ocho
capturas en modo por tiempo del 3 de noviembre, una diferencia de una
sola muestra (\SI{0.08}{\micro\second}). En el barrido por temperatura
del mismo día, ese instante se recorrió de \SI{65.60}{\micro\second} a
\SI{65.84}{\micro\second} mientras la muestra se enfriaba de
\SI{40.4}{\celsius} a \SI{37.9}{\celsius}. En esta sesión el corrimiento
tiene el sentido que corresponde a la menor velocidad del sonido en agua
más fría; en sesiones anteriores se observó el sentido opuesto, una
discrepancia que se discute en la \cref{sec:trabajo_futuro}. No se
cuantifica esa relación, que pertenece a la física de la muestra; lo que
interesa aquí
es que el sistema resuelve y registra un corrimiento de tres muestras
de forma consistente con la temperatura asignada a cada captura.

\subsubsection{Ruido y número de promedios}

El ruido, medido como la desviación estándar de las primeras 400
muestras de la traza, disminuye con el número de adquisiciones
promediadas $N$: de \SI{16.3}{\micro\volt} con $N = 14$ a
\SI{3.4}{\micro\volt} con $N = 396$, un factor de 4.9 frente al 5.3 que
predice la dependencia $1/\sqrt{N}$. Por encima de unas 400
adquisiciones se estanca en un piso de unos \SI{3}{\micro\volt}
(\SIrange{2.8}{3.2}{\micro\volt} con $N$ de 728 a 924), lo que sugiere una
componente que el promedio no elimina; no se identificó su origen.

En el modo por temperatura, $N$ no lo fija el operador sino la duración
de la ventana, es decir, el tiempo que tarda la muestra en cruzar el
margen alrededor de cada objetivo. El ancho de la ventana de integración
se fijó en $\pm\SI{0.02}{\celsius}$ alrededor de cada temperatura
objetivo en la sesión final (en sesiones previas,
$\pm\SI{0.1}{\celsius}$). Este valor no representa la incertidumbre de
la temperatura asignada a cada señal: el DS18B20 resuelve
\SI{0.0625}{\celsius} a 12 bits y su exactitud declarada por el
fabricante es de $\pm\SI{0.5}{\celsius}$; además, en una misma captura
de esa sesión se observaron diferencias de hasta \SI{1.18}{\celsius}
entre los cuatro sensores. La incertidumbre de temperatura está
dominada, por tanto, por la exactitud y la dispersión de los sensores,
no por el ancho de la ventana. Con un margen tan estrecho, $N$ varió de
14 a 140 entre objetivos de la misma sesión, y en dos de ellos la
temperatura registrada al abrir y al cerrar la ventana fue la misma: el
sensor no resolvió el cambio dentro de la ventana. Un margen más
estrecho que la resolución del sensor no mejora la asignación de
temperatura y, en cambio, reduce y vuelve irregular el número de
promedios.

% ---------------------------------------------------------
\subsection{Sesión con un cambio controlado}
\label{sec:sesion_referencia}
% ---------------------------------------------------------

La sesión \texttt{20260825\_192627} fue la primera realizada con el
módulo de temperatura ya integrado a la secuencia automática. Es
anterior a la versión congelada y a su formato de registro, pero es la
única que contiene un cambio deliberado de un parámetro del láser.
Consta de diecisiete capturas de \num{50000} muestras cada una,
repartidas en dos series separadas por una interrupción durante la cual
se reprogramó el microcontrolador y se modificó el retardo del
conmutador electro-óptico del láser.

La primera serie (\texttt{m0002}--\texttt{m0006}) se tomó con un
\textit{EO delay} de \SI{3800}{\micro\second} y con la lectura de
temperatura activa, que descendió de \SI{37.27}{\celsius} a
\SI{36.97}{\celsius}. La segunda (\texttt{m0008}--\texttt{m0017}) se
tomó con un \textit{EO delay} de \SI{3700}{\micro\second} y sin lectura
de temperatura, por la pérdida del módulo que se documenta en la
\cref{tab:fallas_entorno}. Los dos valores del retardo no constan en el
registro de la sesión: su formato, de dieciocho columnas, es anterior a
la incorporación de los parámetros del láser a cada fila.

\subsubsection{Reconstrucción temporal entre escalas distintas}
\label{sec:reconstruccion_escalas}

Las dos series se adquirieron con bases de tiempo diferentes, de modo
que el máximo de la señal aparece en índices de muestra muy distintos
según la serie: alrededor de la muestra \num{13641} en la primera y de
la \num{6729} en la segunda. Reconstruidos los ejes con sus respectivos
valores de \texttt{XINCR}, \texttt{XZERO} y \texttt{PT\_OFF} mediante la
\cref{eq:conv_tiempo}, las diecisiete capturas sitúan el máximo de la
amplitud absoluta entre \SI{69.06}{\micro\second} y
\SI{69.24}{\micro\second}, con media de \SI{69.16}{\micro\second} y
desviación estándar de \SI{0.05}{\micro\second}.

Dos muestreos con resoluciones distintas, aplicados a la misma señal
física, convergen al mismo instante tras la reconstrucción. Eso valida
simultáneamente tres cosas: que el preámbulo se consulta de forma
atómica con la transferencia de la forma de onda, que sus valores se
escriben correctamente en el registro de sesión, y que la fórmula de
conversión los interpreta como corresponde. La
\cref{sec:fallas_entorno} documenta un caso en el que ese emparejamiento
entre muestras y metadatos falló y produjo una reconstrucción falsa que
no se anunciaba como tal.

\subsubsection{Discriminación del cambio}
\label{sec:discriminacion}

Dentro de cada serie, el coeficiente de variación del valor pico a pico
fue de \SI{0.33}{\percent} en la primera y de \SI{0.64}{\percent} en la
segunda, cifras del mismo orden que las del régimen estable de las
sesiones congeladas. La amplitud media pasa de \SI{9.87}{\milli\volt}
en la primera serie a \SI{9.07}{\milli\volt} en la segunda, una
diferencia cercana al \SI{8}{\percent}: más de diez veces la
dispersión interna de cualquiera de las dos.

Tres salvedades acotan lo que puede afirmarse. La primera es que entre
series se perdió también la lectura del módulo de temperatura, de modo
que lo que se controló fue un parámetro y lo que cambió fueron dos; esa
segunda diferencia es de registro y no de excitación, porque la
temperatura no interviene en el cálculo de la amplitud. La segunda es
que la comparación vale solo dentro de la sesión: las sesiones
congeladas muestran derivas de hasta \SI{12.8}{\percent} al inicio de
una secuencia y de hasta \SI{42}{\percent} entre sesiones, de modo que
un cambio del \SI{8}{\percent} entre días distintos no sería
distinguible. La tercera es que la atribución de la diferencia a su
causa descansa en esta sesión sobre información externa al archivo,
porque el formato de entonces no guardaba el retardo del conmutador. El
formato vigente registra en cada fila \texttt{output\_level},
\texttt{eo\_delay\_us} y \texttt{burst\_mode}
(\cref{sec:formato_datos}); con él, un cambio como este queda atribuible
a su causa sin depender de que alguien lo haya anotado aparte.

No se afirma una relación cuantitativa entre el retardo del conmutador
electro-óptico y la amplitud fotoacústica: el arreglo no mide la
energía del pulso, y dos puntos no sostienen una curva. Lo que se afirma
es que el sistema distingue el cambio y conserva su causa.

% ---------------------------------------------------------
\subsection{Alcance de la prueba}
\label{sec:alcance_prueba}
% ---------------------------------------------------------

Ambos modos automáticos se recorrieron de punta a punta con los tres
instrumentos reales y con la versión congelada. Quedan cuatro límites
que conviene declarar.

El primero es la normalización. Sin una medida de la energía de cada
pulso, la deriva al inicio de cada secuencia y la deriva entre sesiones
no pueden atribuirse al láser ni a la muestra, y las amplitudes solo son
comparables dentro de una sesión.

El segundo es el marcador temporal. El máximo de la amplitud absoluta
fue estable en las sesiones del 25 de agosto y del 3 de noviembre, pero
en varias capturas del 5 de noviembre saltó a otros picos de la traza,
a \num{160}, \num{309} y \num{88} microsegundos, cuando cambió la
amplitud relativa entre ellos. Como instante de arribo, el máximo
absoluto no es un criterio robusto, y la identificación de cada pico con
un trayecto de la celda (\cref{ap:tiempos_arribo}) queda como trabajo a
futuro.

El tercero es la temperatura. La incertidumbre de la temperatura
asignada a cada captura la fijan la exactitud y la dispersión de los
sensores, no el margen de la ventana, y el modo por temperatura recorre
solo objetivos en descenso (\cref{sec:secuencia}).

El cuarto es el alcance experimental. Las sesiones cubren agua
destilada entre unos 29 y \SI{42}{\celsius} en enfriamiento pasivo; no
queda establecido el comportamiento del sistema con otros líquidos,
intervalos más amplios o regímenes distintos del enfriamiento pasivo.
